\documentclass[11pt,a4paper]{article}

\usepackage[utf8]{inputenc}
\usepackage[T1]{fontenc}
\usepackage[brazilian]{babel}
\usepackage[top=2cm,bottom=2cm,left=2cm,right=2cm]{geometry}
\usepackage{enumitem}
\usepackage{xcolor}
\usepackage{hyperref}

\definecolor{mainblue}{RGB}{0,74,128}

\hypersetup{
  colorlinks=true,
  linkcolor=mainblue,
  urlcolor=mainblue
}

\setlist[itemize]{leftmargin=*, itemsep=2pt, topsep=3pt}

\title{\textbf{Plano da Apresentação: Ciclo Hamiltoniano}}
\author{}
\date{}

\begin{document}

\maketitle
\vspace{-2em}

\section{Ideia central}

O Ciclo Hamiltoniano é um problema de grafos simples de entender:

\begin{quote}
Dá para sair de um vértice, passar por todos os vértices exatamente uma vez e voltar para o início?
\end{quote}

O ponto central da apresentação é mostrar o contraste entre três ideias:

\begin{itemize}
  \item explicar o problema é fácil;
  \item conferir uma resposta também é fácil;
  \item descobrir se existe uma resposta, em geral, pode ser muito difícil.
\end{itemize}

\section{Objetivos para a turma}

Ao final da apresentação, a turma deve entender:

\begin{itemize}
  \item o que é um caminho hamiltoniano;
  \item o que é um ciclo hamiltoniano;
  \item qual é a diferença entre ciclo hamiltoniano e ciclo euleriano;
  \item como escrever o problema como um problema de decisão;
  \item por que uma solução candidata é fácil de verificar;
  \item por que o problema está em NP;
  \item por que ele é um exemplo clássico de problema NP-completo;
  \item que existem algoritmos exatos, mas eles crescem muito rápido;
  \item que, na prática, muitas abordagens usam busca, poda, programação dinâmica ou heurísticas.
\end{itemize}

\section{Organização sugerida dos slides}

\subsection{Slide 1 -- Título e pergunta central}

Abrir com o título \textbf{Ciclo Hamiltoniano} e uma pergunta simples:

\begin{quote}
Existe um passeio fechado que visita todos os vértices uma única vez?
\end{quote}

Esse começo coloca a turma direto no problema antes de entrar nas definições. Se já houver um grafo pequeno na tela, a pergunta fica concreta desde o primeiro momento.

\subsection{Slide 2 -- Definições básicas}

Alinhar o vocabulário mínimo:

\begin{itemize}
  \item grafo;
  \item vértice;
  \item aresta;
  \item caminho;
  \item ciclo;
  \item caminho hamiltoniano;
  \item ciclo hamiltoniano.
\end{itemize}

Não vale gastar muito tempo aqui. O objetivo é apenas garantir que todos consigam acompanhar os exemplos seguintes.

\subsection{Slide 3 -- Exemplo com ciclo hamiltoniano}

Mostrar um grafo pequeno em que o ciclo existe e escrever a ordem dos vértices. O ponto é deixar claro que não basta o desenho parecer circular: o ciclo precisa visitar todos os vértices uma única vez e voltar ao início.

Esse slide é importante porque a definição fica mais fácil quando a turma vê uma sequência concreta.

\subsection{Slide 4 -- Exemplo sem ciclo hamiltoniano}

Trazer um grafo pequeno que não tem ciclo, de preferência com um vértice de grau 1 ou alguma estrutura que force repetição.

A ideia é mostrar que a ausência do ciclo não é uma questão de desenhar melhor o caminho. Em alguns casos, a própria estrutura do grafo impede a solução.

\subsection{Slide 5 -- Hamiltoniano vs. Euleriano}

Usar a comparação com Euler para evitar uma confusão comum:

\begin{itemize}
  \item o ciclo hamiltoniano se preocupa com vértices;
  \item o ciclo euleriano se preocupa com arestas.
\end{itemize}

A tabela pode ser simples e visual. O objetivo é deixar claro que os dois problemas parecem parecidos no nome, mas pedem coisas diferentes.

\subsection{Slide 6 -- Problema de decisão}

Apresentar a pergunta formal:

\begin{quote}
Dado um grafo $G$, $G$ possui um ciclo hamiltoniano?
\end{quote}

Depois, mostrar três formas de olhar para o tema:

\begin{itemize}
  \item decidir se existe;
  \item encontrar um ciclo quando existe;
  \item otimizar algum custo, como no Problema do Caixeiro Viajante.
\end{itemize}

Esse slide faz a ponte entre o desenho do grafo e a linguagem de complexidade.

\subsection{Slide 7 -- Verificar é fácil}

Mostrar uma sequência candidata, por exemplo:

\begin{quote}
\texttt{A, C, D, B, E, A}
\end{quote}

Explicar como conferir:

\begin{itemize}
  \item todos os vértices aparecem uma vez;
  \item cada par consecutivo é uma aresta;
  \item o último vértice volta ao primeiro.
\end{itemize}

O ponto é separar duas tarefas diferentes: verificar uma resposta pronta é rápido; descobrir a resposta pode não ser.

\subsection{Slide 8 -- Encontrar é difícil}

Usar este slide para falar da explosão combinatória. Com muitos vértices, o número de ordens possíveis cresce muito rápido, então testar tudo deixa de ser viável.

Vale citar que existem algoritmos melhores que força bruta, como programação dinâmica no estilo Held-Karp, mas que eles ainda têm custo exponencial.

\subsection{Slide 9 -- NP e NP-completo sem pesar}

Explicar sem entrar em prova longa:

\begin{itemize}
  \item o problema está em NP porque dá para conferir um certificado em tempo polinomial;
  \item o problema é NP-completo porque representa uma dificuldade central dentro de NP.
\end{itemize}

A ideia não é demonstrar a redução, mas deixar a turma entender o peso da afirmação: se houvesse um algoritmo polinomial geral para Ciclo Hamiltoniano, isso mexeria com a questão $P$ versus $NP$.

\subsection{Slide 10 -- Quando dá para garantir ou tentar resolver}

Mostrar que ``difícil em geral'' não significa ``impossível sempre''.

Existem casos especiais e condições suficientes, como o Teorema de Dirac, que garantem ciclo hamiltoniano em grafos bem conectados.

Na prática, também aparecem abordagens como:

\begin{itemize}
  \item busca com poda;
  \item programação dinâmica;
  \item modelagem como SAT;
  \item programação linear inteira;
  \item heurísticas.
\end{itemize}

Esse slide ajuda a não deixar a apresentação pessimista demais.

\subsection{Slide 11 -- Aplicações e conexões}

Falar de conexões sem vender o problema como uma solução direta para tudo. Boas pontes:

\begin{itemize}
  \item roteamento;
  \item Problema do Caixeiro Viajante;
  \item passeio do cavalo;
  \item reconstrução em bioinformática;
  \item modelos de visita sem repetição em redes.
\end{itemize}

O valor deste slide é mostrar que o problema aparece tanto como ferramenta de modelagem quanto como exemplo clássico de limite algorítmico.

\subsection{Slide 12 -- Fechamento}

Fechar voltando à frase-guia:

\begin{quote}
Simples de explicar, fácil de verificar e difícil de resolver em geral.
\end{quote}

Uma boa pergunta final é:

\begin{quote}
Por que verificar uma resposta não é a mesma coisa que encontrar uma resposta?
\end{quote}

Ela resume a ligação entre o problema de grafos e a discussão de complexidade.

\section{Tópicos para o handout}

O handout pode ser um pouco mais completo que os slides, mas ainda curto.

Sugestão de estrutura:

\begin{itemize}
  \item definição informal;
  \item definição mais formal;
  \item exemplo positivo;
  \item exemplo negativo;
  \item diferença para ciclo euleriano;
  \item problema de decisão;
  \item certificado e verificação;
  \item NP e NP-completude;
  \item ideia de algoritmos: força bruta, backtracking, programação dinâmica e heurísticas;
  \item aplicações e conexões.
\end{itemize}

\end{document}
